\begin{abstract}
We study the strategic interaction between a financial regulator, responsible for balancing market growth and investor protection, and an issuing firm preparing for an initial public offering (IPO). The regulator prescribes an enforcement tolerance level $\varepsilon$, which bounds the budget by which the firm can window-dress its fundamental feature vector into a presented feature vector while jointly choosing the presentation and share price. The market consists of $k \ge 2$ heterogeneous investor subpopulations with distinct personal valuation functions over the presented features. While the firm seeks to maximize total capital raised, the regulator aims to maximize social welfare by facilitating capital raising by viable firms while curbing misallocations to non-viable firms. We model this interaction as a Stackelberg screening game with the regulator as the leader and the firm as the follower.

We analyze this game under two distinct information structures. Under the first, the \emph{Decision Problem}, the feature-valuation distributions are known, and we characterize the firm's optimal response for any enforcement level $\varepsilon$ by reducing its presentation and pricing problem to $2^k-1$ constrained convex optimization subproblems. We further show that expected social welfare need not be monotone in regulatory tolerance: an overly strict regime ($\varepsilon = 0$) can stifle capital raising by viable firms with baseline features associated with low perceived valuations, whereas an overly lax regime ($\varepsilon = \infty$) can enable non-viable firms to mislead investors. Thus, an optimal policy can occur at a strictly intermediate enforcement level. Under the second, the \emph{Learning Problem}, the market distribution is unknown, and the regulator must learn the optimal enforcement policy through repeated market interactions over a finite horizon $T$. Using the established Gaussian Process Thompson Sampling (GP-TS) algorithm with a Matérn-$5/2$ kernel, we obtain a sublinear cumulative regret of $\widetilde{\mathcal{O}}(T^{2/3})$.
\end{abstract}



